% Part: first-order-logic
% Chapter: models-theories
% Section: introduction

\documentclass[../../../include/open-logic-section]{subfiles}

\begin{document}

\olfileid{fol}{mat}{int}
\olsection{Inleiding}

\begin{explain}
De ontwikkeling van de axiomatische methode---een vakgebied opbouwen
vanuit basisbegrippen en axioma's---was een grote stap in de
geschiedenis van de wetenschap, en vooral in die van de wiskunde. Als
je een vak zo opbouwt, moet je eerst een aantal vragen helder krijgen.
Waar gaat het vak over? Met welke begrippen begin je? Hoe hangen die
begrippen samen? Kun je alle andere begrippen daarmee definiëren? En
aan welke wetten houden die begrippen zich, of moeten ze zich houden?

De axiomatische methode en de logica passen precies bij elkaar. Met
formele logica kun je een axiomatische theorie nauwkeurig opschrijven.
Je kunt volgens precies vastgelegde regels bewijzen welke stellingen
uit de axioma's volgen. En je kunt alle systemen die aan die axioma's
voldoen op een vaste, geordende manier onderzoeken.
\end{explain}

\begin{defn}
Een verzameling !!{sentence}s~$\Gamma$ heet \emph{gesloten} precies
dan als elk logisch gevolg van die verzameling er al in zit. In
symbolen: als $\Gamma \Entails !A$, dan $!A \in \Gamma$. De
\emph{afsluiting} van een verzameling !!{sentence}s~$\Gamma$ krijg je
door alle logische gevolgen ervan te nemen. Dat is
$\Setabs{!A}{\Gamma \Entails !A}$.

We zeggen dat~$\Gamma$ \emph{door een verzameling zinnen~$\Delta$
wordt geaxiomatiseerd} als $\Gamma$ precies de afsluiting van~$\Delta$
is. Dan zitten alle zinnen die logisch uit die axioma's volgen er al
in.
\end{defn}

\begin{explain}
Een axiomatische theorie kun je zien als de verzameling van alle zinnen
die logisch uit haar axioma's~$\Delta$ volgen. Anders gezegd: eerst
kies je een eerste-ordetaal. De niet-logische tekens daarin staan voor
de basisbegrippen waarmee het vak begint (de primitieve begrippen).
Daarnaast geef je een verzameling !!{sentence}s die de basiswetten van
het vak uitdrukken. De theorie wordt dan weergegeven door alle
!!{sentence}s in die taal die logisch uit de axioma's volgen. Dat kan
een eenvoudige theorie zijn met maar één basisbegrip en eenvoudige
axioma's, zoals de theorie van partiële ordeningen, maar ook een
ingewikkelde theorie zoals de Newtoniaanse mechanica.

De volgende logische feiten laten zien waarom deze formele aanpak van
de axiomatische methode zo belangrijk is. Neem aan dat~$\Gamma$ het
axiomasysteem van een theorie is. Dat axiomasysteem is dus een
verzameling zinnen.
\begin{enumerate}
\item Met een axiomasysteem kunnen we precies aangeven welke klasse
  !!{structure}s we bedoelen. Stel dat we een bepaalde klasse
  !!{structure}s willen beschrijven. Het axiomasysteem~$\Gamma$
  beschrijft precies die klasse als de !!{structure}s in de klasse
  precies alle~$\Struct M$ zijn waarvoor $\Sat{M}{\Gamma}$.
\item Dat kan misgaan. Er kan een~$\Struct M$ zijn waarvoor
  $\Sat{M}{\Gamma}$, terwijl~$\Struct M$ niet een van de
  !!{structure}s is die we bedoelden. Om zo'n~$\Struct M$ uit te
  sluiten, kunnen we axioma's toevoegen die daarin niet waar zijn.
\item Stel dat~$\Gamma$ in elk geval waar is in alle bedoelde
  !!{structure}s. Dan is een zin~$!A$ in al die !!{structure}s waar
  zodra $\Gamma \Entails !A$. Daarom kunnen we logische hulpmiddelen
  gebruiken, zoals methoden voor formele !!{derivation}. Om te laten
  zien dat een zin in alle bedoelde !!{structure}s waar is, hoeven we
  dan alleen te laten zien dat die zin logisch uit de axioma's volgt.
\item Soms beginnen we niet met een bepaalde klasse bedoelde
  !!{structure}s, maar met de axioma's zelf. We kiezen een paar
  basisbegrippen en leggen in axioma's vast aan welke wetten ze moeten
  voldoen. Eerst willen we weten of alle axioma's tegelijk waar kunnen
  zijn. Als ze elkaar tegenspreken, kun je de basisbegrippen niet zo
  uitleggen dat alle wetten gelden; de theorie hangt dan niet samen.
  Je kunt laten zien dat dit probleem zich niet voordoet door één model
  van~$\Gamma$ te vinden. Als de theorie modellen heeft, kunnen we die
  met logische methoden onderzoeken en ook nieuwe modellen bouwen.
\item Ook willen we weten of elk axioma echt nodig is. Dat is de vraag
  naar de onafhankelijkheid van de axioma's. Misschien volgt één
  axioma al uit de andere axioma's. Dat axioma is dan overbodig
  (redundant). Een axioma~$!A$ in~$\Gamma$ is overbodig als we kunnen
  bewijzen dat $\Gamma \setminus \{!A\} \Entails !A$. Om te bewijzen
  dat het niet overbodig is, laten we zien dat
  $(\Gamma \setminus \{!A\}) \cup \{\lnot !A\}$ vervulbaar is: alle
  zinnen in die verzameling kunnen dus tegelijk waar zijn. Zo is
  aangetoond dat het parallellenaxioma onafhankelijk is van de andere
  axioma's van de meetkunde.
\item Ten slotte vragen we welke begrippen je met andere begrippen kunt
  definiëren. De taal die je kiest, bepaalt wat de modellen van een
  theorie bevatten. Maar niet elk onderdeel van een theorie hoeft
  apart in die modellen te staan. Iedere ordening~$\le$ bepaalt
  bijvoorbeeld een bijbehorende strikte ordening~$<$. Als je de ene
  hebt, kun je de andere definiëren. Een model met zo'n ordening hoeft
  de bijbehorende strikte ordening dus niet \emph{ook} als aparte
  relatie te bevatten. De algemene vraag is: wanneer kun je een
  relatie met andere relaties definiëren? En wanneer kan dat niet,
  zodat je die relatie als een van de basisbegrippen (primitieve
  begrippen) van de taal moet opnemen?
\end{enumerate}
\end{explain}


\end{document}
